\subsection{Описание метода}

\subsubsection{Определение квадратурных формул}

Рассмотрим задачу численного вычисления определённого интеграла от функции $f(x)$ на отрезке $[a, b]$:
\[
I = \int_a^b f(x) \, dx.
\]

\textbf{Квадратурной формулой} \textit{(или аппроксимирующей квадратурой)} называется выражение вида
\[
Q_n[f] = \sum_{i=0}^{n} w_i f(x_i),
\]
где $\{x_i\}_{i=0}^n \subset [a, b]$ --- узлы квадратуры, а $\{w_i\}_{i=0}^n$ --- весовые коэффициенты. 
Квадратурная формула аппроксимирует интеграл $I$ конечной линейной комбинацией значений подынтегральной функции. 

\textit{Погрешность квадратуры} определяется как разность следующего вида:
\[
E_n[f] = I - Q_n[f]
\]

\subsubsection{Степень точности квадратурной формулы}

\textbf{Алгебраической степенью точности} квадратурной формулы называется наибольшее целое неотрицательное число $m$, такое что формула точна для всех алгебраических многочленов степени не выше $m$:
\[
\forall k=\overline{0,m}\ E_n[x^k] = 0, \quad \text{но} \quad E_n[x^{m+1}] \neq 0
\]

Иными словами, если подынтегральная функция является многочленом $P_m(x)$ степени $m$, то квадратурная формула даёт точное значение интеграла. 

\subsubsection{Формула трапеций}

Рассмотрим отрезок \( [a,b] \) и построим интерполяционный многочлен Лагранжа первой степени на узлах \( \{a,b\} \):
\[
  L_1(x)=\frac{x-b}{a-b}f(a)+\frac{x-a}{b-a}f(b)
\]

Определим остаточный член \( R_1(x) \) согласно теореме об остаточном члене интерполяционной формулы Лагранжа:
\[
  f(x) = L_1(x) + R_1(x), \quad R_1(x) = \frac{f''(\xi(x))}{2}(x-a)(x-b), \quad \xi(x) \in (a,b)
\]

Проинтегрируем обе части равенства по отрезку \( [a,b] \):
\[
  \int_a^b f(x)\,dx = \int_a^b L_1(x)\,dx + \int_a^b R_1(x)\,dx
\]

Вычислим интеграл от интерполяционного многочлена почленно:
\[
  \begin{aligned}
    \int_a^b L_1(x)\,dx &= f(a)\int_a^b \frac{x-b}{a-b}\,dx + f(b)\int_a^b \frac{x-a}{b-a}\,dx \\
    &= \frac{f(a)}{a-b} \left[ \frac{(x-b)^2}{2} \right]_a^b + \frac{f(b)}{b-a} \left[ \frac{(x-a)^2}{2} \right]_a^b \\
    &= \frac{f(a)}{a-b} \left( 0 - \frac{(a-b)^2}{2} \right) + \frac{f(b)}{b-a} \left( \frac{(b-a)^2}{2} - 0 \right) \\
    &= \frac{b-a}{2}f(a) + \frac{b-a}{2}f(b) = \frac{b-a}{2}\bigl(f(a)+f(b)\bigr)
  \end{aligned}
\]

Для оценки погрешности проинтегрируем остаточный член. Поскольку \((x-a)(x-b) \le 0\) на \([a,b]\), а \(f''\) непрерывна, по интегральной теореме о среднем:
\[
  \int_a^b \dfrac{f''(\xi(x))}{2}(x-a)(x-b)\,dx = \frac{f''(\xi)}{2} \int_a^b (x-a)(x-b)\,dx, \quad \xi \in (a,b)
\]

Вычислим интеграл от квадратного многочлена:
\[
  \int_a^b (x-a)(x-b)\,dx = \int_a^b \left(x^2 - (a+b)x + ab\right)dx = -\frac{(b-a)^3}{6}
\]

Подставляя результат, получаем остаточный член формулы трапеций:
\[
  \int_a^b R_1(x)\,dx = -\frac{(b-a)^3}{12} f''(\xi)
\]

Окончательно, полная формула трапеций записывается в виде:
\[
  \int_a^b f(x)\,dx = \frac{f(a)+f(b)}{2}\cdot h - \frac{h^3}{12} f''(\xi), \quad \xi \in (a,b), \quad h=b-a
\]

Можно заметить, что алгебраическая степень точности метода трапеций равна $m = 1$. Это следует из погрешности квадратуры: если \( f \) --- многочлен первой степени, то \( R_1(x)\equiv 0 \).

\subsubsection{Обработка точек разрыва}

Если функция $f(x)$ имеет разрывы на отрезке интегрирования, квадратурная формула может давать некорректный результат или приводить к исключительным ситуациям \textit{(деление на ноль, переполнение)}.

\begin{enumerate}
  \item \textbf{Устранимый разрыв или скачок.} Интеграл Римана существует, и алгоритм должен заменить значение $f(c)$ на одно из значений $f(c-\varepsilon(h)),\ f(c+\varepsilon(h))$, где $\varepsilon(h)$ --- некоторая малая величина.

  \item \textbf{Бесконечный разрыв.} Интеграл Римана не существует, и алгоритм должен прервать вычисления, уведомив пользователя об ошибке.
\end{enumerate}

Таким образом, в реализации метода трапеций мы должны проанализировать подынтегральную функцию на наличие особенностей и адаптировать разбиение интервала интегрирования в зависимости от типа разрыва.